\chapter{Spacetime Uniqueness}
\label{ch:spacetime}

The $\Dz$ operator uniquely determines a 4-dimensional Lorentzian spacetime.
The AF-channel gives 1 temporal dimension ($\re(\afcoeff) = 0$, $\im(\afcoeff) > 0$),
the SY-channel gives 3 spatial dimensions ($\phiS$ rank 3),
and the signature $(1,3)$ follows from
the non-compact/compact duality $\vp\ps = -1$ vs.\ $\zs\zsc = +1$.

\section{\texorpdfstring{$\so(3,1)$}{so(3,1)}: Lorentz Algebra}

\begin{definition}[$I_4$]
\label{def:I4}
\lean{FUST.Physics.Lorentz.I4}
$I_4 = \mathrm{Fin}\,1 \oplus \mathrm{Fin}\,3$: spacetime index.
\end{definition}

\begin{theorem}[$\so(3,1)$ dimension]
\label{thm:so31}
\lean{FUST.Physics.Lorentz.finrank_so31}
$\so(3,1) \cong \R^6$ with $\dim\so(3,1) = 6$:
3 rotations (spatial antisymmetric block) and 3 boosts (space-time symmetric block).
\end{theorem}

\section{\texorpdfstring{Poincar\'{e} Algebra $\iso(3,1)$}{Poincaré Algebra iso(3,1)}}

\begin{theorem}[Translation dimension]
\label{thm:finrank-translations}
\lean{FUST.Physics.Poincare.finrank_translations}
$\dim(I_4 \to \R) = 4$.
\end{theorem}

\begin{theorem}[Poincar\'{e} dimension]
\label{thm:finrank-poincare}
\lean{FUST.Physics.Poincare.finrank_poincare}
$\dim\iso(3,1) = \dim\so(3,1) + \dim\R^4 = 6 + 4 = 10$.
\end{theorem}

\begin{definition}[Minkowski bilinear form]
\label{def:minkowski}
\lean{FUST.Physics.Poincare.minkowskiBilin}
$\eta_{\mu\nu} = \mathrm{diag}(+1, -1, -1, -1)$ on $I_4$.
\end{definition}

\section{Spacetime from \texorpdfstring{$\Dz$}{Dζ}}

\begin{theorem}[Spacetime uniqueness]
\label{thm:spacetime-uniqueness}
\lean{FUST.SpacetimeUniqueness.spacetime_from_Dζ}
The $\Dz$ operator determines $I_4 = \mathrm{Fin}\,1 \oplus \mathrm{Fin}\,3$:
\begin{enumerate}
  \item $\re(\afcoeff) = 0$, $\im(\afcoeff) > 0$: 1 temporal channel ($\phiA$),
  \item $\phiS$ rank $= 3$: 3 spatial channels (Diff5, Diff3, Diff2),
  \item $\phiA[\mathrm{id}](1) \ne 0$: temporal channel nontrivial,
  \item Signature $(1,3)$: $(6a)^2 > 0$ (temporal), $(\afcoeff \cdot b)^2 < 0$ (spatial),
  \item Lorentz invariance: $\eta(\Lambda v, \Lambda w) = \eta(v,w)$
    for all $\Lambda \in O(3,1)$,
  \item $\vp$-dilation $=$ translation in log-coordinates:
    $\exp(t + \log\vp) = \vp \cdot \exp(t)$,
  \item 4 linearly independent translation vectors from $\vp$-scaling on $I_4$.
\end{enumerate}
\end{theorem}

\begin{theorem}[Gauge--spacetime independence]
\label{thm:gauge-spacetime-indep}
\lean{FUST.SpacetimeUniqueness.gauge_spacetime_independence}
Gauge symmetry ($\C$-linearity of $\Fz$, star $\ne$ id)
and spacetime symmetry (Lorentz invariance of $\eta$, star $=$ id on $\R$)
act on independent degrees of freedom.
Total: $12 + 10 = 22$ symmetry dimensions.
\end{theorem}

\section{General Relativity}

\begin{definition}[Lorentz connection]
\label{def:lorentz-connection}
\lean{FUST.Physics.Gravity.LorentzConnection}
$\omega : I_4 \to \so(3,1)$, a field of $\so(3,1)$-valued 1-forms.
\end{definition}

\begin{definition}[Curvature]
\label{def:curvature}
\lean{FUST.Physics.Gravity.curvature}
$F_{\mu\nu} = [\omega_\mu, \omega_\nu]$.
\end{definition}

\begin{theorem}[Bianchi identity]
\label{thm:bianchi}
\lean{FUST.Physics.Gravity.bianchi_identity}
$F_{[\mu\nu,\rho]} = 0$.
\end{theorem}

\begin{theorem}[Vacuum Einstein equations]
\label{thm:vacuum-einstein}
\lean{FUST.Physics.Gravity.vacuum_einstein}
$G_{\mu\nu} = R_{\mu\nu} - \tfrac{1}{2}R\,g_{\mu\nu} = 0$.
\end{theorem}
